\documentclass[11pt,a4paper]{article}
\usepackage[margin=20mm]{geometry}
\usepackage{fontspec}
\setmainfont{DejaVu Serif}
\setsansfont{DejaVu Sans}
\setmonofont{DejaVu Sans Mono}[Scale=0.85]
\usepackage{polyglossia}
\setdefaultlanguage{russian}
\usepackage{graphicx}
\usepackage{hyperref}
\usepackage{enumitem}
\usepackage{fancyvrb}
\hypersetup{hidelinks}
\setlength{\parindent}{0pt}
\setlength{\parskip}{4pt}
\usepackage{titlesec}
\titleformat{\section}{\large\bfseries}{}{0pt}{}
\titlespacing*{\section}{0pt}{10pt}{4pt}
\sloppy
\setlist{nosep,leftmargin=*}
\emergencystretch=3em
\newcommand{\shot}[2]{\begin{center}\includegraphics[width=\linewidth,height=.34\textheight,keepaspectratio]{#1}\par\small #2\end{center}}
\begin{document}
\begin{center}
{\small Университет ИТМО · Фронт-энд разработка · 2026}\par
{\LARGE Домашняя работа № 5}\par
{\large npm и основы Vue}\par
Светлов Ярослав, группа J3211
\end{center}
\section*{Цель и технологии}
Освоить npm и компонентный подход Vue на отдельном приложении заметок. Использованы Vue 3, Vite, Vue Router, Pinia, Axios, Bootstrap и JSON Server.
\section*{Компоненты и данные}
\begin{itemize}
\item BaseLayout: общая навигация и slot для содержимого.
\item NoteCard: получает заметку через props и отображает заголовок, текст и дату.
\item NoteForm: реактивная форма v-model, проверка данных, событие saved.
\item NotesPage: загрузка в onMounted, список v-for, состояния v-if. NewNotePage показывает форму и после сохранения возвращает к списку.
\item stores/notes.js: общее состояние Pinia и запросы Axios. Сервер проверяет заголовок и текст и сохраняет данные в рабочую базу.
\end{itemize}

\section*{Команды npm}
npm install устанавливает зависимости и обновляет lock-файл. npm ci воспроизводит установку по существующему lock-файлу. npm run запускает сценарии из package.json. Проверены сборка, тест API и npm audit: на момент проверки аудит показал 0 уязвимостей.
\section*{Проверка}
В Firefox 148 и Chromium проверены список, переход к форме, создание заметки с многострочным текстом и сохранение после обновления. Пробельный заголовок отклоняется. На ширине 390 px переполнения нет. Выключенный API вызывает понятную ошибку и позволяет повторить загрузку. Автотест проверяет пустой список, валидацию, создание, нормализацию заголовка, запись на диск и закрытый служебный маршрут.
\section*{Запуск и проверка}
Требуется Node.js 22.12 или новее. Команды выполняются из папки этой работы.
\begin{verbatim}
npm ci
npm run api
\end{verbatim}
Во втором терминале из той же папки:
\begin{verbatim}
npm start
\end{verbatim}
Открыть \url{http://localhost:8085}. Автотесты: \texttt{npm test}. Сборка: \texttt{npm run build}; просмотр сборки: \texttt{npm run preview} вместо dev-сервера. API должен оставаться запущенным.
 API: \url{http://localhost:3005}. Вход не требуется. Новые заметки сохраняются в \texttt{db.local.json}.\section*{Вывод}
Реализованы маршрутизация, компоненты, общее состояние и серверное сохранение. Полные исходники находятся в src.

\clearpage
\section*{Скриншоты}
\shot{checks/form.png}{Рисунок 1. NoteForm внутри BaseLayout.}
\shot{checks/notes.png}{Рисунок 2. NoteCard и список заметок на узком экране.}
\clearpage\section*{Пример исходного компонента}
Компонент \texttt{src/components/NoteCard.vue} из текущей работы:
\VerbatimInput[fontsize=\footnotesize]{src/components/NoteCard.vue}
Данные приходят через обязательное свойство note. Карточка не выполняет запросы и не хранит отдельную копию состояния. Остальной код расположен в src/components, src/views и src/stores.

\end{document}
